% year: תשפ"ג
% type: מבחן
% semester: א
% moed: א
% number: 4א
% points: 10
% categories: lhopital
% source: exams/מבחנים/תשפג/מועד א תשפג.pdf

\begin{question}
חשבו את הגבול $\displaystyle\lim_{x\to 1}\frac{x^2+\ln(x)-1}{e^x-e}$.
\end{question}

\begin{solution}
בהצבת $x=1$: המונה $1+0-1=0$ והמכנה $e-e=0$, כלומר ביטוי מהצורה $\frac00$. המונה והמכנה גזירים בסביבת $1$, ונגזרת המכנה $e^x\neq 0$. לפי \textbf{כלל לופיטל}:
\[ \lim_{x\to1}\frac{x^2+\ln x-1}{e^x-e}=\lim_{x\to1}\frac{2x+\frac1x}{e^x}=\frac{2+1}{e}=\frac{3}{e}, \]
כאשר הגבול האחרון מחושב בהצבה (רציפות), ולכן גם הגבול המקורי קיים ושווה ל-$\frac{3}{e}$.
\end{solution}
